% ============================================================================
% Chuong 2: CO SO LY THUYET VA CONG NGHE
% ============================================================================

\chapter{CƠ SỞ LÝ THUYẾT VÀ CÔNG NGHỆ}
\label{ch:foundations}

% Logo cong nghe: moi logo MOT hinh rieng co caption, rong >= 1/5 trang.
\newcommand{\techlogo}[3]{% #1=file #2=caption #3=label
  \begin{figure}[H]
    \centering
    \includegraphics[width=0.22\linewidth]{../shared/figures/logos/#1}
    \caption{#2}
    \label{#3}
  \end{figure}}

Chương này trình bày cơ sở lý thuyết của các giải thuật cốt lõi và các công nghệ nền
tảng được sử dụng trong đồ án. Bối cảnh kiến trúc (sự dịch chuyển từ một khối sang vi
dịch vụ), các mẫu tấn công cần phát hiện và lý do chọn quan sát ở tầng nhân đã được
trình bày ở Chương~\ref{ch:overview}; chương này tập trung đi sâu vào \emph{cách} hiện
thực. Bám sát đúng đường đi của một sự kiện trong hệ thống, chương lần lượt làm rõ: cơ
chế thu thập ở tầng nhân bằng eBPF (Mục~\ref{sec:theory-ebpf}), ước lượng tần suất bằng
Count-Min Sketch để phát hiện quét cổng (Mục~\ref{sec:theory-cms}), phân tích đồ thị và
phân loại địa chỉ (Mục~\ref{sec:graph-ip}), và cuối cùng là nền tảng công nghệ triển
khai (Mục~\ref{sec:tech-stack}).

\section{Thu thập sự kiện ở tầng nhân bằng eBPF}
\label{sec:theory-ebpf}

\subsection{eBPF và mô hình thực thi an toàn trong nhân Linux}

\techlogo{ebpf.png}{Logo eBPF.}{fig:logo-ebpf}
eBPF (\textit{extended Berkeley Packet Filter}) là một máy ảo nhỏ chạy ngay bên
trong nhân Linux, cho phép nạp và thực thi các chương trình do người dùng cung
cấp tại nhiều điểm móc (\textit{hook}) của nhân mà không cần biên dịch lại nhân
hay nạp mô-đun nhân \cite{gregg2019bpf}. Điểm mấu chốt khiến eBPF phù hợp cho bài
toán giám sát an ninh là cơ chế \textit{verifier}: trước khi một chương trình
eBPF được nạp, bộ kiểm chứng của nhân phân tích tĩnh toàn bộ đồ thị luồng điều
khiển để bảo đảm chương trình luôn kết thúc (không có vòng lặp vô hạn), không
truy cập bộ nhớ ngoài phạm vi cho phép, và không thực hiện thao tác nguy hiểm.
Sau khi vượt qua bộ kiểm chứng, chương trình được biên dịch tức thời (\textit{JIT})
sang mã máy gốc nên chạy với hiệu năng gần như mã nhân biên dịch sẵn. So với việc
viết một mô-đun nhân truyền thống --- vốn có thể làm sập (\textit{panic}) toàn bộ
hệ điều hành nếu mắc lỗi và phải biên dịch lại theo từng phiên bản nhân --- eBPF
mang lại sự an toàn và khả năng triển khai linh hoạt hơn hẳn, đây chính là lý do
nó trở thành nền tảng cho phần lớn công cụ quan sát an ninh hiện đại trên
Kubernetes. Hình~\ref{fig:ebpf-exec} tóm tắt vòng đời của một chương trình eBPF:
mã C được biên dịch sang bytecode BPF bằng \texttt{clang/LLVM}, nạp vào nhân qua
lời gọi \texttt{bpf()}, đi qua bộ kiểm chứng (nếu không an toàn sẽ bị từ chối ngay),
được JIT sang mã máy gốc, rồi gắn tại điểm móc và tương tác với các BPF map.

\begin{figure}[H]
  \centering
  \resizebox{0.9\linewidth}{!}{% Mo hinh thuc thi eBPF: bien dich -> nap -> verifier -> JIT -> gan hook + maps. Chuong 2.
\begin{tikzpicture}[font=\footnotesize, node distance=0.6cm and 1.2cm,
    box/.style ={draw, rounded corners=2pt, align=center, minimum height=0.85cm,
                 minimum width=2.9cm},
    us/.style  ={box, fill=blue!10},
    kn/.style  ={box, fill=green!12},
    map/.style ={draw, rounded corners=2pt, fill=orange!16, align=center,
                 minimum height=1.0cm, minimum width=2.7cm},
    flow/.style={-{Stealth[length=2mm]}, thick},
  ]
  % --- Khong gian nguoi dung (trai) ---
  \node[us] (src) {Mã nguồn C\\(\texttt{*.bpf.c})};
  \node[us, below=of src] (clang) {clang/LLVM\\$\rightarrow$ bytecode BPF};
  \node[us, below=of clang] (loader) {libbpf loader\\(\texttt{bpf()} syscall)};
  \node[draw, dashed, fit=(src)(clang)(loader), inner sep=8pt,
        label=above:{\textbf{Không gian người dùng}}] (ubox) {};

  % --- Nhan Linux (phai): chuoi doc verifier -> jit -> hook; maps ben phai hook ---
  \node[kn, right=2.6cm of src] (verif) {Verifier\\(kiểm chứng an toàn:\\dừng, giới hạn bộ nhớ)};
  \node[kn, below=of verif] (jit) {JIT compiler\\$\rightarrow$ mã máy gốc};
  \node[kn, below=of jit] (hook) {Chương trình eBPF (đã JIT)\\gắn tại \texttt{kprobe}};
  \node[map, right=1.1cm of hook] (maps) {BPF maps\\(LPM\_TRIE,\\RINGBUF, HASH)};
  \node[draw, dashed, fit=(verif)(jit)(hook)(maps), inner sep=9pt,
        label=above:{\textbf{Nhân Linux}}] (kbox) {};

  % Luong chinh
  \draw[flow] (src) -- (clang);
  \draw[flow] (clang) -- (loader);
  \draw[flow] (loader.east) -- ++(0.45,0) |- (verif.west);   % nap vao verifier
  \draw[flow] (verif) -- node[right, font=\scriptsize]{đạt} (jit);
  \draw[flow] (jit) -- (hook);
  \draw[flow, {Stealth}-{Stealth}] (hook) -- (maps);          % chuong trinh <-> maps
  % Bi loai neu khong an toan (do, dut net) quay ve loader
  \draw[flow, red, dashed] (verif.west) to[bend left=18] ($(loader.east)+(0.25,0.35)$);
  \node[red, font=\scriptsize, align=center] at ($(verif.west)!0.5!(loader.east)+(0,0.62)$)
    {bị loại nếu\\không an toàn};
\end{tikzpicture}
}
  \caption{Mô hình thực thi của một chương trình eBPF.}
  \label{fig:ebpf-exec}
\end{figure}

Một chương trình eBPF trao đổi dữ liệu với không gian người dùng và lưu trạng thái
qua các cấu trúc gọi là \textit{BPF map}. Bộ thu thập của đồ án dùng ba map, mỗi
map một vai trò (Hình~\ref{fig:bpf-maps}). \texttt{cidr\_map} kiểu
\texttt{BPF\_MAP\_TYPE\_LPM\_TRIE} phân loại địa chỉ IP theo dải mạng (CIDR) ngay
trong nhân theo nguyên tắc khớp tiền tố dài nhất (Mục~\ref{sec:graph-ip}), nhờ đó
lọc bỏ sớm những kết nối không thuộc mạng pod, giảm tải cho tầng người dùng.
\texttt{events} kiểu \texttt{BPF\_MAP\_TYPE\_RINGBUF} là bộ đệm vòng đa nhà sản
xuất, một nhà tiêu thụ, dùng để chuyển sự kiện từ nhân lên không gian người dùng;
so với \textit{perf buffer} thế hệ trước (cấp phát bộ đệm theo từng CPU và dễ mất
thứ tự sự kiện giữa các CPU), ring buffer bảo toàn thứ tự toàn cục và dùng bộ nhớ
hiệu quả hơn, phù hợp cho luồng kết nối mật độ cao. Cuối cùng, \texttt{sock\_store}
kiểu \texttt{BPF\_MAP\_TYPE\_HASH} giữ con trỏ \texttt{sock} mà chương trình lưu
lúc vào hàm để đọc lại lúc hàm trả về --- chi tiết kỹ thuật cần thiết sẽ được làm
rõ ngay sau đây.

\begin{figure}[H]
  \centering
  \resizebox{0.98\linewidth}{!}{% Ba BPF map cua collector: cidr_map, events, sock_store. Chuong 2.
\begin{tikzpicture}[font=\footnotesize, node distance=0.7cm,
    m/.style={draw, rounded corners=3pt, align=left, minimum width=4.7cm,
              minimum height=2.0cm, text width=4.4cm, inner sep=6pt},
  ]
  \node[m, fill=orange!14] (lpm)
    {\textbf{\texttt{cidr\_map}} --- \texttt{BPF\_MAP\_TYPE\_LPM\_TRIE}\\[2pt]
     \footnotesize key: \texttt{lpm\_key\{prefixlen, addr\}}\\
     value: \texttt{\_\_u32} (nhãn vùng mạng)\\
     max\_entries: 1024\\[2pt]
     \textit{Phân loại IP theo CIDR ngay trong nhân (khớp tiền tố dài nhất).}};
  \node[m, fill=green!12, right=of lpm] (rb)
    {\textbf{\texttt{events}} --- \texttt{BPF\_MAP\_TYPE\_RINGBUF}\\[2pt]
     \footnotesize dung lượng: $256\times1024$ B\\[2pt]
     \textit{Bộ đệm vòng đa nhà sản xuất, một nhà tiêu thụ; chuyển} \texttt{network\_event\_t} \textit{từ nhân lên user, bảo toàn thứ tự.}};
  \node[m, fill=blue!10, right=of rb] (hs)
    {\textbf{\texttt{sock\_store}} --- \texttt{BPF\_MAP\_TYPE\_HASH}\\[2pt]
     \footnotesize key: \texttt{\_\_u64} (\texttt{pid\_tgid})\\
     value: \texttt{struct sock *}\\
     max\_entries: 4096\\[2pt]
     \textit{Giữ con trỏ} \texttt{sock} \textit{từ kprobe (vào) để kretprobe (ra) đọc lại.}};
\end{tikzpicture}
}
  \caption{Ba BPF map của bộ thu thập và vai trò tương ứng.}
  \label{fig:bpf-maps}
\end{figure}

\subsection{Lựa chọn điểm móc và tính khả chuyển nhờ CO-RE}

Nhân Linux cho phép gắn chương trình eBPF tại nhiều tầng khác nhau, mỗi tầng có
sự đánh đổi riêng. \texttt{XDP} và \texttt{tc} hoạt động rất sớm ở tầng trình điều
khiển mạng, lý tưởng cho việc lọc/loại bỏ gói với thông lượng lớn nhưng chỉ thấy
gói thô, khó truy được ngữ cảnh tiến trình. \textit{Tracepoint} là các điểm móc
tĩnh do nhân định nghĩa sẵn, ổn định nhưng kém linh hoạt. \textit{kprobe} là điểm
móc động cho phép chèn vào gần như bất kỳ hàm nào của nhân. Đồ án chọn gắn
\texttt{kprobe} (và \texttt{kretprobe} tương ứng) vào hàm \texttt{tcp\_v4\_connect}
--- hàm được nhân gọi mỗi khi một tiến trình khởi tạo kết nối TCP/IPv4 qua lời gọi
\texttt{connect()}. Tại điểm này, chương trình đọc được đầy đủ địa chỉ nguồn,
địa chỉ đích, cổng đích và mã tiến trình, vốn là toàn bộ thông tin cần thiết để
phát hiện hành vi quét cổng và dịch chuyển ngang. Mô hình quan sát ở đây mang
tính nền tảng: vì mọi kết nối TCP của mọi pod trên một nút đều phải đi qua nhân
của nút đó, một chương trình eBPF đặt tại \texttt{tcp\_v4\_connect} sẽ quan sát
được trọn vẹn lưu lượng kết nối mà không cần cài đặt tác nhân vào từng container.
Hình~\ref{fig:kprobe-stack} cho thấy vị trí của hai điểm móc trên đường gọi hàm
nhân, còn Bảng~\ref{tab:hook-compare} đối chiếu các lựa chọn điểm móc.

\begin{figure}[H]
  \centering
  \resizebox{0.9\linewidth}{!}{% Vi tri kprobe/kretprobe tren ham tcp_v4_connect trong ngan xep TCP. Chuong 2.
\begin{tikzpicture}[font=\footnotesize, node distance=0.55cm,
    lyr/.style ={draw, rounded corners=2pt, align=center, minimum height=0.8cm,
                 minimum width=4.8cm, fill=gray!10},
    hook/.style={draw, rounded corners=2pt, align=center, fill=green!14,
                 minimum height=0.85cm, minimum width=3.0cm},
    flow/.style={-{Stealth[length=2mm]}, thick},
  ]
  \node[lyr, fill=blue!10] (app) {Tiến trình trong Pod: \texttt{connect(fd, dst, len)}};
  \node[lyr, below=of app] (sys) {Lời gọi hệ thống \texttt{\_\_sys\_connect}};
  \node[lyr, below=of sys, fill=yellow!16] (fn) {\texttt{tcp\_v4\_connect(sk, uaddr, len)}};
  \node[lyr, below=of fn] (hs) {Thiết lập kết nối / bắt tay 3 bước};

  % kprobe + kretprobe gan tren tcp_v4_connect; tach xa nhau theo chieu doc
  \node[hook, right=2.4cm of fn, yshift=1.05cm] (kentry) {\texttt{kprobe} (vào):\\lưu con trỏ \texttt{sock}};
  \node[hook, right=2.4cm of fn, yshift=-1.05cm] (kret) {\texttt{kretprobe} (ra):\\đọc \texttt{sock}, sinh sự kiện};

  \draw[flow] (app) -- (sys);
  \draw[flow] (sys) -- (fn);
  \draw[flow] (fn) -- (hs);
  \draw[{Stealth}-, thick, green!50!black] (kentry.west) -- ($(fn.east)+(0,0.18)$);
  \draw[{Stealth}-, thick, green!50!black] (kret.west) -- ($(fn.east)+(0,-0.18)$);
  \node[font=\scriptsize\itshape, align=center, above=0.18cm of kentry]
    {gắn động vào cùng một hàm nhân};
\end{tikzpicture}
}
  \caption{Vị trí gắn \texttt{kprobe}/\texttt{kretprobe} trong ngăn xếp TCP.}
  \label{fig:kprobe-stack}
\end{figure}

\begin{table}[H]
  \centering
  \caption{So sánh các điểm móc eBPF cho bài toán quan sát kết nối.}
  \label{tab:hook-compare}
  \small
  \begin{tabularx}{\linewidth}{@{}l l X@{}}
    \toprule
    \textbf{Điểm móc} & \textbf{Tầng} & \textbf{Đặc điểm / phù hợp} \\
    \midrule
    XDP / tc        & Trình điều khiển / hàng đợi & Rất sớm, thông lượng cao, nhưng chỉ thấy gói thô, khó truy ngữ cảnh tiến trình \\
    tracepoint      & Điểm tĩnh của nhân          & Ổn định theo phiên bản nhưng cố định, kém linh hoạt \\
    \textbf{kprobe} & \textbf{Hàm nhân (động)}    & \textbf{Gắn được vào \texttt{tcp\_v4\_connect}; đọc đủ src/dst/port/pid --- lựa chọn của đồ án} \\
    \bottomrule
  \end{tabularx}
\end{table}

Cơ chế thu thập được tách làm hai nửa vì lý do kỹ thuật: tại thời điểm \emph{vào}
hàm (\texttt{kprobe}), địa chỉ đích chưa chắc đã được điền đầy đủ vào cấu trúc
\texttt{sock}, nên chương trình chỉ lưu con trỏ \texttt{sock} theo khóa
\texttt{pid\_tgid} vào \texttt{sock\_store}; tại thời điểm \emph{trả về}
(\texttt{kretprobe}), chương trình đọc lại con trỏ này, trích ra địa chỉ nguồn,
địa chỉ đích và cổng đích, tra \texttt{cidr\_map} để xác định nhãn vùng mạng, và
chỉ khi nguồn thuộc mạng pod (\texttt{src\_label} $= 1$) mới cấp một ô trong
\texttt{events} (ring buffer) để điền cấu trúc \texttt{network\_event\_t} gồm
\{\texttt{src\_ip}, \texttt{dst\_ip}, \texttt{dst\_port}, \texttt{pid},
\texttt{timestamp\_ns}\} rồi đẩy lên không gian người dùng. Bộ lọc theo nhãn ngay
trong nhân này loại bỏ phần lớn lưu lượng không liên quan tại nguồn, giảm tải đáng
kể cho tầng phân tích. Hình~\ref{fig:packet-path} mô tả toàn bộ đường đi của một
sự kiện từ lời gọi \texttt{connect()} trong pod tới khi tới được bộ máy pipeline.

\begin{figure}[H]
  \centering
  \resizebox{0.6\linewidth}{!}{% Duong di chi tiet: connect() -> kprobe/kretprobe -> loc LPM -> RINGBUF -> collector -> pipeline. Chuong 2.
\begin{tikzpicture}[font=\footnotesize, node distance=0.62cm and 0.9cm,
    pod/.style ={draw, rounded corners=2pt, fill=blue!10, align=center, minimum height=0.8cm, minimum width=4.4cm},
    ker/.style ={draw, rounded corners=2pt, fill=green!12, align=center, minimum height=0.8cm, minimum width=5.6cm},
    dec/.style ={draw, diamond, aspect=2.6, fill=yellow!20, align=center, inner sep=1pt},
    usr/.style ={draw, rounded corners=2pt, fill=blue!8, align=center, minimum height=0.8cm, minimum width=5.6cm},
    map/.style ={draw, rounded corners=2pt, fill=orange!16, align=center, minimum height=0.7cm, font=\scriptsize},
    flow/.style={-{Stealth[length=2mm]}, thick},
  ]
  \node[pod] (a) {\textbf{Pod:} \texttt{connect(dst\_ip:dst\_port)}};
  \node[ker, below=of a] (b) {\textbf{kprobe} \texttt{tcp\_v4\_connect\_entry}:\\lưu con trỏ \texttt{sock} theo \texttt{pid\_tgid}};
  \node[ker, below=of b] (c) {\textbf{kretprobe} \texttt{tcp\_v4\_connect\_return}:\\đọc \texttt{sock} $\rightarrow$ trích \texttt{src\_ip, dst\_ip, dst\_port}};
  \node[ker, below=of c] (d) {Tra \texttt{cidr\_map} (LPM-Trie)\\$\rightarrow$ \texttt{src\_label}, \texttt{dst\_label}};
  \node[dec, below=of d] (e) {\texttt{src\_label}\\$=1$ (pod)?};
  \node[ker, below=1.0cm of e] (f) {Cấp ô \texttt{events} (RINGBUF), điền\\\texttt{network\_event\_t}\{src,dst,port,pid,ts\} $\rightarrow$ submit};
  \node[usr, below=of f] (g) {\textbf{Collector} (libbpf, \texttt{ring\_buffer\_\_poll})\\$\rightarrow$ dòng JSON $\rightarrow$ FIFO $\rightarrow$ \textbf{pipeline}};
  \node[draw, rounded corners=2pt, fill=gray!12, right=of e, align=center, minimum height=0.7cm] (drop) {bỏ qua\\(không phải pod)};

  % maps on the side
  \node[map, right=0.7cm of b] (m1) {\texttt{sock\_store}\\HASH};
  \node[map, right=0.7cm of d] (m2) {\texttt{cidr\_map}\\LPM\_TRIE};
  \node[map, right=0.7cm of f] (m3) {\texttt{events}\\RINGBUF};

  \draw[flow] (a)--(b); \draw[flow] (b)--(c); \draw[flow] (c)--(d); \draw[flow] (d)--(e);
  \draw[flow] (e)-- node[right,font=\scriptsize]{đúng} (f);
  \draw[flow] (e)-- node[above,font=\scriptsize]{sai} (drop);
  \draw[flow] (f)--(g);
  \draw[{Stealth}-{Stealth}, gray] (b.east)--(m1.west);
  \draw[{Stealth}-{Stealth}, gray] (d.east)--(m2.west);
  \draw[{Stealth}-{Stealth}, gray] (f.east)--(m3.west);
\end{tikzpicture}
}
  \caption{Đường đi của một sự kiện kết nối qua eBPF.}
  \label{fig:packet-path}
\end{figure}

Để một chương trình eBPF biên dịch một lần có thể chạy trên nhiều phiên bản nhân
khác nhau, đồ án dựa vào cơ chế CO-RE (\textit{Compile Once -- Run Everywhere})
kết hợp với thông tin kiểu BTF (\textit{BPF Type Format}) mà nhân hiện đại đính
kèm. Thay vì giả định bố cục bộ nhớ của các cấu trúc nhân tại thời điểm biên dịch
(điều buộc phải có header trùng khít với từng nút), chương trình truy cập trường
dữ liệu thông qua các phép định vị lại (\textit{relocation}) được giải tại thời
điểm nạp dựa trên BTF của nhân đích. Nhờ đó cùng một tệp đối tượng eBPF chạy được
trên cụm có nhân không đồng nhất, giảm đáng kể chi phí vận hành. Cuối cùng, cần
ghi nhận một giới hạn vốn có của lựa chọn điểm móc này: vì \texttt{kprobe} trên
\texttt{tcp\_v4\_connect} chỉ kích hoạt khi có lời gọi \texttt{connect()}, các
kiểu quét gửi gói SYN thô (ví dụ \texttt{nmap -sS} chạy với quyền root) sẽ không
đi qua đường này và do đó không bị quan sát; phân tích định lượng về giới hạn này
được trình bày ở Chương~\ref{ch:eval}.

\section{Ước lượng tần suất trên luồng: Count-Min Sketch}
\label{sec:theory-cms}

Bài toán cốt lõi mà tầng phát hiện phải giải là: trên một luồng sự kiện kết nối
gần như vô hạn và tốc độ cao, ước lượng \emph{tần suất xuất hiện} của mỗi địa chỉ
nguồn để nhận ra nguồn nào đang tạo ra số lượng kết nối bất thường. Lưu trữ chính
xác tần suất bằng bảng băm đòi hỏi bộ nhớ tỉ lệ thuận với số phần tử phân biệt ---
không khả thi khi số địa chỉ và cổng có thể rất lớn và biến động liên tục. Count-Min
Sketch (CMS) \cite{cormode2005cms} giải quyết vấn đề này bằng một cấu trúc xác suất
sử dụng bộ nhớ dưới tuyến tính, đánh đổi một sai số ước lượng có thể kiểm soát được
để lấy bộ nhớ cố định và thời gian cập nhật/truy vấn hằng số.

Về cấu trúc, CMS là một ma trận đếm kích thước $d \times w$ đi kèm $d$ hàm băm độc
lập đôi một, mỗi hàm ánh xạ một khóa vào một trong $w$ cột (đồ án dùng họ hàm băm
phi mật mã MurmurHash \cite{appleby2011murmur} vì tốc độ cao và phân bố tốt). Thao
tác \texttt{record(x)} tăng một đơn vị tại $d$ ô $(\,i, h_i(x)\,)$ với mọi hàng $i$;
thao tác \texttt{estimate(x)} trả về \emph{giá trị nhỏ nhất} trong $d$ ô đó. Tính
chất nền tảng của cấu trúc là nó \emph{không bao giờ ước lượng thiếu}: do va chạm
băm chỉ có thể làm các ô bị cộng thêm, ta luôn có $\hat{f}(x) \ge f(x)$. Mặt khác,
sai số vượt lên trên bị chặn theo xác suất: với chiều rộng $w = \lceil e/\varepsilon \rceil$
và chiều sâu $d = \lceil \ln(1/\delta) \rceil$, ước lượng thỏa mãn
$\hat{f}(x) \le f(x) + \varepsilon \lVert \mathbf{f} \rVert_1$ với xác suất ít nhất
$1-\delta$. Phát biểu và chứng minh đầy đủ của cận này, cùng phân tích so sánh với
Misra--Gries và Count Sketch, đã được trình bày trong Báo cáo học phần CNPM~1 nên ở
đây chỉ nêu lại kết quả phục vụ phần thiết kế. Hình~\ref{fig:cms} minh họa trực quan
cấu trúc, còn Giải thuật~\ref{alg:cms} tóm tắt hai thao tác cơ bản.

\begin{figure}[H]
  \centering
  \resizebox{0.8\linewidth}{!}{% Cau truc Count-Min Sketch: ma tran d x w + bank ham bam, truy van = min. Chuong 2.
% Thiet ke tranh giao cat: x -> cot hash box -> o duoc to dam (mui ten ngang, khong cat nhan).
\begin{tikzpicture}[font=\footnotesize, >={Stealth[length=1.7mm]}, x=0.66cm, y=0.66cm,
    cell/.style={draw, minimum size=0.62cm, inner sep=0},
    hb/.style={draw, rounded corners=1pt, fill=green!12, minimum width=0.9cm,
               minimum height=0.5cm, font=\scriptsize, inner sep=1pt},
  ]
  % Luoi d=5 hang x w=8 cot (cot ve tu x=2..9 de chua bank ham bam ben trai)
  \foreach \r in {1,...,5}{
    \foreach \c in {1,...,8}{ \node[cell] at (\c+1,-\r) {}; }
  }
  % O ma khoa x roi vao (1 o moi hang)
  \foreach \r/\c in {1/3,2/6,3/2,4/7,5/4}{ \node[cell, fill=blue!22] at (\c+1,-\r) {}; }
  % Bank ham bam: cot cac hop h1..h5 thang hang voi tung hang luoi
  \foreach \r in {1,...,5}{ \node[hb] (h\r) at (-0.1,-\r) {$h_{\r}$}; }
  % Khoa x
  \node[draw, rounded corners=2pt, fill=orange!22, font=\small] (x) at (-2.7,-3) {khóa $x$};
  % x -> cac hop ham bam (fan tren vung trong, khong cat nhan)
  \foreach \r in {1,...,5}{ \draw[->, gray!75] (x.east) -- (h\r.west); }
  % moi hop ham bam -> o to dam tuong ung (mui ten NGANG, sach)
  \foreach \r/\c in {1/3,2/6,3/2,4/7,5/4}{ \draw[->, blue!65] (h\r.east) -- ({\c+1-0.32},-\r); }
  % nhan w (tren) va d (trai)
  \draw[decorate, decoration={brace, amplitude=4pt}] (1.68,-0.45) -- (9.32,-0.45)
    node[midway, above=3pt] {$w$ cột (chiều rộng)};
  \node[rotate=90, font=\scriptsize, anchor=south] at (-3.9,-3) {$d$ hàng băm};
  % truy van min (ben phai, cach luoi)
  \node[draw, rounded corners=2pt, fill=gray!8, align=left, anchor=west] at (10.3,-3)
    {$\hat f(x)=\displaystyle\min_{1\le i\le d} C[i][\,h_i(x)\,]$\\[2pt]
     lấy giá trị \emph{nhỏ nhất} trong các ô tô đậm\\
     $\Rightarrow$ không bao giờ ước lượng thiếu};
\end{tikzpicture}
}
  \caption{Cấu trúc Count-Min Sketch.}
  \label{fig:cms}
\end{figure}

\begin{algorithm}[H]
\caption{Count-Min Sketch --- cập nhật và truy vấn}
\label{alg:cms}
\begin{algorithmic}[1]
\Function{Record}{$x$}
  \For{$i \gets 1$ \textbf{to} $d$}
    \State $C[i]\big[\,h_i(x) \bmod w\,\big] \gets C[i]\big[\,h_i(x) \bmod w\,\big] + 1$
  \EndFor
\EndFunction
\Function{Estimate}{$x$}
  \State \Return $\min_{1 \le i \le d} \; C[i]\big[\,h_i(x) \bmod w\,\big]$
\EndFunction
\end{algorithmic}
\end{algorithm}

Để thấy rõ cơ chế, Bảng~\ref{tab:cms-trace} minh hoạ trạng thái ma trận sau một
luồng ngắn chạy tay.

\begin{table}[H]
  \centering
  \caption{Trạng thái ma trận CMS ($d=2$, $w=4$) sau luồng $S=(a,b,a,c,a,b)$, với
           $h_1{:}\;a{\to}1, b{\to}3, c{\to}1$ và $h_2{:}\;a{\to}2, b{\to}0, c{\to}3$.}
  \label{tab:cms-trace}
  \small
  \begin{tabular}{@{}l cccc@{}}
    \toprule
    & cột 0 & cột 1 & cột 2 & cột 3 \\
    \midrule
    hàng $h_1$ & 0 & 4 & 0 & 2 \\
    hàng $h_2$ & 2 & 0 & 3 & 1 \\
    \bottomrule
  \end{tabular}
\end{table}

\noindent Khi đó $\hat f(a)=\min\big(C[1][1], C[2][2]\big)=\min(4,3)=3=f(a)$ (chính
xác), và $\hat f(c)=\min\big(C[1][1], C[2][3]\big)=\min(4,1)=1=f(c)$: dù ô $C[1][1]$
bị va chạm giữa $a$ và $c$ (giá trị 4), phép lấy nhỏ nhất vẫn cho ước lượng đúng nhờ
hàng còn lại --- đây chính là trực giác đằng sau tính chất không-ước-lượng-thiếu.

Trong hệ thống của đồ án, CMS được dùng với khóa là \emph{địa chỉ IP nguồn}: mỗi sự
kiện kết nối làm tăng bộ đếm của địa chỉ nguồn tương ứng, và khi giá trị ước lượng
của một nguồn vượt một ngưỡng định trước trong một cửa sổ thời gian, hệ thống phát
một cảnh báo \texttt{port\_scan}. Cấu hình thực tế dùng ma trận $5 \times 2048$ với
ngưỡng $40$ trên cửa sổ trượt kiểu \textit{tumbling} dài $60$ giây. Cần nhấn mạnh
một đặc điểm thiết kế có hệ quả đến độ chính xác: bộ đếm khóa theo địa chỉ nguồn nên
nó đếm \emph{tổng lưu lượng kết nối} của nguồn đó, chứ không đếm \emph{số cổng đích
phân biệt}. Lựa chọn này khiến bộ phát hiện thực chất là một bộ phát hiện theo tốc
độ kết nối; ưu, nhược điểm và các hệ quả định lượng của nó (bỏ sót khi quét chậm,
báo giả khi một client hợp lệ tạo nhiều kết nối) được phân tích kỹ ở
Chương~\ref{ch:eval}.

\section{Phân tích đồ thị và phân loại địa chỉ trong mạng nội cụm}
\label{sec:graph-ip}
\label{ch:other-dsa}

Ngoài việc đếm tần suất từng nguồn riêng lẻ, hệ thống còn dựng một \emph{đồ thị
truyền thông} trong đó mỗi đỉnh là một địa chỉ và mỗi cung biểu diễn một kết nối đã
quan sát được. Trên đồ thị này, hai bài toán đồ thị kinh điển được vận dụng để phát
hiện những mẫu hành vi mà phép đếm tần suất đơn lẻ không nhìn thấy, kèm theo một bài
toán phân loại địa chỉ phục vụ cho cả tầng nhân lẫn tầng phân tích.

\subsection{Tarjan --- phát hiện thành phần liên thông mạnh}

Một thành phần liên thông mạnh (\textit{Strongly Connected Component}, SCC) của đồ
thị có hướng là một tập đỉnh cực đại trong đó mọi cặp đỉnh đều có thể đến được nhau.
Trong ngữ cảnh an ninh nội cụm, một SCC không tầm thường (kích thước lớn hơn một)
giữa các pod là một dấu hiệu đáng ngờ: nó cho thấy một nhóm pod đang liên lạc vòng
tròn lẫn nhau, mẫu hình thường gặp ở các mạng điều khiển (\textit{command-and-control})
hoặc ở giai đoạn dịch chuyển ngang khi kẻ tấn công đã kiểm soát nhiều điểm. Đồ án sử
dụng giải thuật Tarjan \cite{tarjan1972depth} để liệt kê toàn bộ SCC trong một lần
duyệt theo chiều sâu (DFS): mỗi đỉnh được gán một chỉ số khám phá và một giá trị
\textit{low-link} --- chỉ số nhỏ nhất có thể quay về được qua các cung ngược --- đồng
thời được đẩy vào một ngăn xếp; khi một đỉnh có low-link bằng chính chỉ số khám phá
của nó, toàn bộ các đỉnh phía trên nó trong ngăn xếp tạo thành một SCC. Giải thuật
chạy trong thời gian tuyến tính $O(V+E)$ theo số đỉnh và số cung, đủ nhanh để chạy
định kỳ trên đồ thị truyền thông đang lớn dần của cụm. Giải thuật~\ref{alg:tarjan}
trình bày ý tưởng cốt lõi của một lần duyệt.

\begin{algorithm}[H]
\caption{Tarjan --- liệt kê thành phần liên thông mạnh}
\label{alg:tarjan}
\begin{algorithmic}[1]
\State $index \gets 0$;\quad ngăn xếp $S \gets \varnothing$
\Function{StrongConnect}{$v$}
  \State $v.index \gets index$;\; $v.low \gets index$;\; $index \gets index + 1$
  \State đẩy $v$ vào $S$;\; $v.onStack \gets \mathbf{true}$
  \ForAll{cung $(v, w)$}
    \If{$w$ chưa được thăm}
      \State \Call{StrongConnect}{$w$};\; $v.low \gets \min(v.low,\, w.low)$
    \ElsIf{$w.onStack$}
      \State $v.low \gets \min(v.low,\, w.index)$
    \EndIf
  \EndFor
  \If{$v.low = v.index$}
    \State lấy các đỉnh khỏi $S$ cho tới $v$ $\;\Rightarrow\;$ một thành phần liên thông mạnh
  \EndIf
\EndFunction
\end{algorithmic}
\end{algorithm}

\begin{figure}[H]
  \centering
  \resizebox{0.62\linewidth}{!}{% Vi du do thi truyen thong: mot SCC kha nghi {A,B,C} + cac dinh khac. Chuong 2.
\begin{tikzpicture}[font=\footnotesize, >={Stealth[length=2mm]},
    v/.style={draw, circle, minimum size=0.85cm, fill=blue!8},
    s/.style={draw, circle, minimum size=0.85cm, fill=red!14},
  ]
  \node[s] (A) at (0,0) {A};
  \node[s] (B) at (1.6,0.7) {B};
  \node[s] (C) at (1.6,-0.7) {C};
  \node[v] (D) at (-1.8,0) {D};
  \node[v] (E) at (3.4,0) {E};
  \node[v] (F) at (3.4,-1.4) {F};
  % SCC cycle A->B->C->A
  \draw[->, thick] (A) -- (B);
  \draw[->, thick] (B) -- (C);
  \draw[->, thick] (C) -- (A);
  % outside edges
  \draw[->] (D) -- (A);
  \draw[->] (B) -- (E);
  \draw[->] (C) -- (F);
  % highlight SCC
  \node[draw, dashed, red, ellipse, fit=(A)(B)(C), inner sep=4pt,
        label={[red,font=\scriptsize]below:{SCC khả nghi (liên thông mạnh)}}] {};
\end{tikzpicture}
}
  \caption{Ví dụ thành phần liên thông mạnh (SCC) trong đồ thị truyền thông.}
  \label{fig:tarjan-ex}
\end{figure}

\subsection{BFS --- bán kính ảnh hưởng theo $k$ bước}

Khi một pod bị nghi ngờ đã bị chiếm, một câu hỏi vận hành quan trọng là: từ pod đó,
kẻ tấn công có thể với tới những đâu trong một số bước giới hạn? Hình thức hóa, ta cần
tính tập khả-đạt-$k$-bước $R_k(s) = \{\, v : \text{tồn tại đường đi từ } s \text{ tới }
v \text{ với độ dài} \le k \,\}$. Vì đồ thị truyền thông không trọng số và ta quan
tâm tới khoảng cách theo số bước, duyệt theo chiều rộng (BFS) là lựa chọn tự nhiên và
tối ưu: BFS mở rộng đồ thị theo từng lớp khoảng cách tăng dần, nên chỉ cần dừng ở lớp
thứ $k$ là thu được toàn bộ tập khả-đạt mong muốn trong thời gian $O(V+E)$
\cite{cormen2009algorithms}. Kích thước của $R_k(s)$ chính là \emph{bán kính ảnh
hưởng} (\textit{blast radius}) của nguồn $s$ --- một chỉ số trực quan giúp người vận
hành ưu tiên xử lý: một nguồn vừa vượt ngưỡng quét cổng lại vừa có bán kính ảnh hưởng
lớn là sự cố cần can thiệp khẩn cấp. Hệ thống sử dụng $k = 3$ bước trong cấu hình mặc
định; Giải thuật~\ref{alg:bfs} mô tả phép duyệt theo lớp có giới hạn độ sâu.

\begin{algorithm}[H]
\caption{BFS giới hạn $k$ bước --- bán kính ảnh hưởng $R_k(s)$}
\label{alg:bfs}
\begin{algorithmic}[1]
\State $R \gets \{s\}$;\quad hàng đợi $Q \gets \big[(s, 0)\big]$
\While{$Q \neq \varnothing$}
  \State $(u, d) \gets$ lấy đầu hàng đợi $Q$
  \If{$d = k$} \State \textbf{bỏ qua}, xét phần tử kế tiếp \EndIf
  \ForAll{láng giềng $v$ của $u$}
    \If{$v \notin R$}
      \State $R \gets R \cup \{v\}$;\; thêm $(v, d{+}1)$ vào $Q$
    \EndIf
  \EndFor
\EndWhile
\State \Return $R$
\end{algorithmic}
\end{algorithm}

\begin{figure}[H]
  \centering
  \resizebox{0.72\linewidth}{!}{% BFS theo lop k-hop tu nguon s; ban kinh anh huong R_k. Chuong 2.
\begin{tikzpicture}[font=\footnotesize, >={Stealth[length=2mm]},
    v/.style={draw, circle, minimum size=0.8cm, fill=blue!10},
    s/.style={draw, circle, minimum size=0.8cm, fill=orange!22},
  ]
  \node[s] (s) at (0,0) {s};
  \node[v] (a) at (2.2,0.9) {a};
  \node[v] (b) at (2.2,-0.9) {b};
  \node[v] (c) at (4.4,1.3) {c};
  \node[v] (d) at (4.4,0) {d};
  \node[v] (e) at (6.6,0.6) {e};
  \draw[->] (s)--(a); \draw[->] (s)--(b);
  \draw[->] (a)--(c); \draw[->] (a)--(d); \draw[->] (b)--(d);
  \draw[->] (d)--(e);
  % layer separators + labels
  \foreach \x/\lab in {0/{lớp 0}, 2.2/{lớp 1}, 4.4/{lớp 2}, 6.6/{lớp 3}}{
    \draw[gray!40, dashed] (\x,-1.7) -- (\x,1.9);
    \node[gray, font=\scriptsize] at (\x,2.1) {\lab};
  }
  \node[align=center, font=\scriptsize] at (3.3,-2.2)
    {$R_3(s)=\{s,a,b,c,d,e\}$ --- bán kính ảnh hưởng trong $\le 3$ bước};
\end{tikzpicture}
}
  \caption{Bán kính ảnh hưởng theo BFS $k$ bước từ nguồn $s$.}
  \label{fig:bfs-ex}
\end{figure}

\subsection{LPM Trie --- phân loại địa chỉ theo tiền tố dài nhất}

Cả tầng nhân lẫn tầng phân tích đều cần trả lời nhanh câu hỏi ``địa chỉ này thuộc
vùng mạng nào'' để phân biệt lưu lượng nội bộ pod, lưu lượng dịch vụ, lưu lượng nút
và lưu lượng bên ngoài. Đây chính là bài toán khớp tiền tố dài nhất
(\textit{Longest Prefix Match}): cho một tập dải mạng CIDR có gán nhãn, hãy tìm dải
\emph{cụ thể nhất} (tiền tố dài nhất) chứa một địa chỉ cho trước. LPM Trie giải bài
toán này bằng một cây nhị phân theo bit của tiền tố: mỗi dải mạng được chèn dọc theo
đường đi tương ứng với các bit tiền tố của nó kèm nhãn, còn thao tác tra cứu đi từ
gốc theo các bit của địa chỉ và ghi nhớ nhãn tại nút khớp sâu nhất. So với cách băm
toàn bộ địa chỉ (không biểu diễn được quan hệ bao hàm giữa các dải) hay cây Radix
tổng quát, LPM Trie cho thời gian tra cứu chặn theo độ dài địa chỉ và biểu diễn tự
nhiên thứ bậc bao hàm của các dải. Trong đồ án, cùng một bảng phân loại được dùng ở
hai nơi: trong nhân qua \texttt{BPF\_MAP\_TYPE\_LPM\_TRIE} để lọc sớm, và trong
pipeline qua bộ phân loại \texttt{ip\_classifier}. Bảng phân loại gán nhãn cho các
dải tiêu biểu: \texttt{10.244.0.0/16} là mạng pod, \texttt{10.96.0.0/12} là mạng
dịch vụ, \texttt{192.168.0.0/16} là mạng nút, và \texttt{0.0.0.0/0} là lưu lượng bên
ngoài; nhờ đó mỗi sự kiện đều mang ngữ cảnh vùng mạng phục vụ cho việc lọc và phân
tích về sau. Giải thuật~\ref{alg:lpm} mô tả thao tác tra cứu nhãn.

\begin{algorithm}[H]
\caption{LPM Trie --- tra cứu nhãn theo tiền tố dài nhất}
\label{alg:lpm}
\begin{algorithmic}[1]
\Function{Lookup}{địa chỉ $a$}
  \State $node \gets root$;\quad $label \gets \textsc{None}$
  \ForAll{bit $b$ của $a$, từ cao xuống thấp}
    \If{$node$ có gán nhãn} \State $label \gets node.label$ \EndIf
    \If{$node.child[b]$ rỗng} \State \textbf{dừng vòng lặp} \EndIf
    \State $node \gets node.child[b]$
  \EndFor
  \State \Return $label$ \Comment{nhãn của tiền tố khớp sâu nhất}
\EndFunction
\end{algorithmic}
\end{algorithm}

\begin{figure}[H]
  \centering
  \resizebox{0.96\linewidth}{!}{% Vi du LPM-Trie phan loai CIDR + tra cuu tien to dai nhat. Chuong 2.
% Duong tra cuu duoc TO DAM; hop giai thich dat DUOI cay, tach khoi moi nut.
\begin{tikzpicture}[font=\footnotesize, >={Stealth[length=2mm]},
    n/.style  ={draw, rounded corners=2pt, align=center, minimum height=0.72cm,
                fill=gray!10, inner sep=3pt},
    pre/.style={draw, rounded corners=2pt, align=center, minimum height=0.72cm,
                fill=green!16, inner sep=3pt},
    hl/.style ={blue!70, line width=1.4pt},
  ]
  % Cay tien to (trie) gan nhan vung mang
  \node[n]   (r)   at (0,0)      {gốc \texttt{0.0.0.0/0}\\$\rightarrow$ external (0)};
  \node[n]   (a)   at (-3.0,-2.0){\texttt{10.0.0.0/8}\\(mạng riêng)};
  \node[pre] (nod) at (3.0,-2.0) {\texttt{192.168.0.0/16}\\$\rightarrow$ node (3)};
  \node[pre] (svc) at (-4.8,-4.1){\texttt{10.96.0.0/12}\\$\rightarrow$ service (2)};
  \node[pre] (pod) at (-1.2,-4.1){\texttt{10.244.0.0/16}\\$\rightarrow$ pod (1)};
  % Cac canh; duong tra cuu r -> a -> pod duoc to dam
  \draw[hl,->] (r) -- (a);
  \draw[->]    (r) -- (nod);
  \draw[->]    (a) -- (svc);
  \draw[hl,->] (a) -- (pod);
  \node[blue!70, font=\scriptsize] at (-3.05,-3.05) {đường tra cứu};
  % Hop giai thich dat DUOI cung, can giua, khong canh nao cat qua
  \node[draw, dashed, rounded corners, align=center, fill=blue!4, inner sep=6pt] (q) at (0.2,-5.9)
    {\textbf{Tra cứu} \texttt{10.244.0.5}: duyệt theo từng bit tiền tố; nút khớp \emph{sâu nhất}
     $=$ \texttt{10.244.0.0/16} $\Rightarrow$ nhãn \textbf{pod (1)} (tiền tố dài nhất thắng).};
\end{tikzpicture}
}
  \caption{Ví dụ phân loại CIDR bằng LPM-Trie.}
  \label{fig:lpm-ex}
\end{figure}

Bốn giải thuật trên không hoạt động rời rạc mà phối hợp trên cùng một dòng xử lý sự
kiện. Hình~\ref{fig:detect-flow} tóm tắt cách chúng kết hợp: mỗi sự kiện kết nối lần
lượt được phân loại theo vùng địa chỉ (LPM Trie), được đếm tần suất theo nguồn (CMS)
để kiểm tra ngưỡng quét cổng, đồng thời được dùng để bồi đắp đồ thị truyền thông ---
nơi giải thuật Tarjan và phép duyệt BFS chạy để phát hiện chu trình khả nghi và ước
lượng bán kính ảnh hưởng. Cách bố trí này cho phép phát hiện trực tuyến với chi phí
thấp ngay khi sự kiện vừa được sinh ra.

\begin{figure}[H]
  \centering
  \resizebox{0.62\linewidth}{!}{% Luu do xu ly mot su kien ket noi qua cac giai thuat. Dung o Chuong 2.
\begin{tikzpicture}[font=\footnotesize, node distance=0.75cm and 1.4cm,
    start/.style={draw, rounded corners=9pt, fill=gray!12, align=center,
                  minimum height=0.8cm, minimum width=3cm},
    proc/.style ={draw, fill=blue!10, align=center, minimum height=0.8cm,
                  minimum width=3.4cm},
    dec/.style  ={draw, diamond, aspect=2.2, fill=yellow!20, align=center,
                  inner sep=1pt},
    alert/.style={draw, fill=orange!20, align=center, minimum height=0.8cm,
                  minimum width=3.6cm},
    flow/.style ={-{Stealth[length=2mm]}, thick},
  ]
  \node[start] (e) {Sự kiện \texttt{connect()}};
  \node[proc, below=of e] (lpm) {Phân loại IP (LPM Trie)};
  \node[proc, below=of lpm] (cms) {CMS: \texttt{record(src\_ip)}};
  \node[dec, below=of cms] (chk) {ước lượng\\$> \theta$ ?};
  \node[alert, right=of chk] (al) {Cảnh báo \texttt{port\_scan}\\$+$ bán kính ảnh hưởng (BFS)};
  \node[proc, below=1.2cm of chk] (graph) {Cập nhật đồ thị truyền thông};
  \node[proc, below=of graph] (scc) {Định kỳ: Tarjan SCC\\(phát hiện chu trình khả nghi)};

  \draw[flow] (e) -- (lpm);
  \draw[flow] (lpm) -- (cms);
  \draw[flow] (cms) -- (chk);
  \draw[flow] (chk) -- node[above, font=\scriptsize]{đúng} (al);
  \draw[flow] (chk) -- node[right, font=\scriptsize]{sai} (graph);
  \draw[flow] (al) |- (graph);
  \draw[flow] (graph) -- (scc);
\end{tikzpicture}
}
  \caption{Lưu đồ xử lý một sự kiện kết nối qua bốn giải thuật cốt lõi.}
  \label{fig:detect-flow}
\end{figure}

\section{Nền tảng công nghệ sử dụng}
\label{sec:tech-stack}

Một đặc trưng của đồ án là toàn bộ hệ thống được xây dựng trên một nền tảng
(\textit{platform}) điện toán đám mây gốc (\textit{cloud-native}) thay vì chạy trực
tiếp trên máy chủ vật lý theo cách truyền thống. Việc lựa chọn nền tảng này không chỉ
là quyết định kỹ thuật phụ trợ mà còn quy định cách thức ứng dụng đặc thù của đồ án
vận hành, do đó mục này trình bày chi tiết từng tầng nền tảng và kết lại bằng một sơ
đồ tổng thể (Hình~\ref{fig:platform}) làm rõ điểm khác biệt so với các kiến trúc đi
trước. Bảng~\ref{tab:tech} tóm tắt công dụng cụ thể của từng công cụ trong hệ thống.

\begin{table}[H]
  \centering
  \caption{Các công cụ/công nghệ nền tảng và công dụng trong hệ thống.}
  \label{tab:tech}
  \small
  \begin{tabularx}{\linewidth}{@{}l X@{}}
    \toprule
    \textbf{Công cụ} & \textbf{Công dụng trong hệ thống} \\
    \midrule
    Kubernetes              & Điều phối container: lập lịch, mở rộng, tự phục hồi; \texttt{DaemonSet} đặt bộ thu thập trên mọi nút \\
    CNI Calico              & Mạng pod và \emph{thực thi} \texttt{NetworkPolicy} --- nền cho vòng phản hồi cách ly \\
    containerd              & Runtime chạy container \\
    Nhân Linux 6.8 + BTF    & Cung cấp eBPF và thông tin kiểu BTF cho CO-RE \\
    eBPF / libbpf           & Thu thập sự kiện kết nối ở tầng nhân (kprobe \texttt{tcp\_v4\_connect}) \\
    NATS JetStream          & Trục thông điệp xuất bản--đăng ký giữa pipeline và các dịch vụ SOC \\
    PostgreSQL              & Lưu trữ bền vững sự cố, cảnh báo, hành động, nhật ký kiểm toán (ACID, JSONB) \\
    Prometheus / Grafana    & Thu thập số liệu và trực quan hóa cho người vận hành \\
    ArgoCD (GitOps)         & Đồng bộ trạng thái cụm về đúng khai báo trong Git (prune, self-heal) \\
    Tekton                  & Tích hợp liên tục (CI) bằng pipeline chạy ngay trong cụm: xây, kiểm thử và đóng gói ảnh OCI \\
    Harbor                  & Registry ảnh container riêng: lưu trữ, quét lỗ hổng và ký ảnh trước khi triển khai \\
    clang/LLVM, bpftool     & Biên dịch chương trình eBPF, sinh \texttt{vmlinux.h}, thao tác BPF \\
    \bottomrule
  \end{tabularx}
\end{table}

Phần dưới đây trình bày từng công nghệ theo cùng một mạch: \emph{nó là gì}, \emph{đồ án
dùng nó làm gì}, và \emph{vì sao chọn nó thay cho các giải pháp thay thế}. Các tầng nhân
Linux, eBPF/libbpf và clang/LLVM/bpftool đã được phân tích sâu ở Mục~\ref{sec:theory-ebpf}
nên không nhắc lại ở đây.

\subsection{Kubernetes --- nền điều phối container}

\techlogo{kubernetes.png}{Logo Kubernetes.}{fig:logo-k8s}
Kubernetes là hệ điều phối (\textit{orchestrator}) container mã nguồn mở, khởi nguồn từ
Google và hiện do CNCF quản trị; nó tự động hóa việc lập lịch, mở rộng, tự phục hồi và
quản lý vòng đời của các container trên một cụm nhiều nút. Đơn vị triển khai nhỏ nhất là
\textit{pod}, và mọi trạng thái mong muốn đều được mô tả \emph{khai báo} bằng tệp YAML.

Trong đồ án, Kubernetes là nền chạy toàn bộ hệ thống: nó điều phối các dịch vụ SOC và
pipeline phân tích như những \textit{Deployment} co giãn, đồng thời --- và đây là điểm
mấu chốt --- dùng đối tượng \texttt{DaemonSet} để bảo đảm bộ thu thập eBPF chạy đúng
\emph{một bản trên mỗi nút}, tự động theo cụm khi thêm/bớt nút. Nhờ vậy độ phủ quan sát
tầng nhân là toàn cục mà không cần thao tác thủ công.

Sở dĩ chọn Kubernetes thay cho các công cụ điều phối nhẹ hơn như Docker Compose hay Docker
Swarm là vì hai lý do đặc thù của bài toán: thứ nhất, chỉ Kubernetes có nguyên hình
\texttt{DaemonSet} đúng ngữ nghĩa ``một tiến trình nền trên mỗi nút'' mà mô hình thu thập
yêu cầu; thứ hai, Kubernetes định nghĩa chuẩn \texttt{NetworkPolicy} --- nền tảng để khép
kín vòng phản hồi cách ly ở Chương~\ref{ch:design}. Đây cũng là chuẩn thực tế của ngành,
giúp đồ án bám sát môi trường triển khai thật.

\subsection{CNI Calico --- mạng pod và thực thi NetworkPolicy}

\techlogo{calico.png}{Logo Project Calico.}{fig:logo-calico}
Calico là một trình cắm mạng (\textit{Container Network Interface}, CNI) cung cấp khả năng
kết nối mạng giữa các pod và, quan trọng hơn, \emph{thực thi} chính sách mạng
\texttt{NetworkPolicy}. Khác với một số CNI chỉ lo định tuyến gói, Calico hiện thực một bộ
máy chính sách ở mỗi nút để cho phép hoặc chặn lưu lượng theo nhãn, không gian tên và dải
địa chỉ.

Vai trò của Calico trong đồ án nằm ở \emph{tầng phản hồi}: khi một nguồn bị xác định là
đang quét cổng, hệ thống sinh ra một \texttt{NetworkPolicy} kiểu chặn lối ra
(\textit{egress-deny}) và Calico là thành phần thực sự áp đặt việc chặn đó trong mạng nội
cụm. Nói cách khác, Calico biến một quyết định phát hiện thành một hành động ngăn chặn cụ
thể --- đúng tinh thần ``phát hiện đi đôi với phòng ngừa'' của tên đề tài.

Calico được chọn vì nó hỗ trợ đầy đủ và ổn định \texttt{NetworkPolicy} (kể cả chính sách
egress), trong khi một số CNI phổ biến như Flannel không thực thi chính sách mạng. So với
Cilium --- một CNI dựa trên eBPF cũng rất mạnh --- Calico cho mô hình chính sách rõ ràng,
tài liệu trưởng thành và đủ cho phạm vi đồ án, đồng thời tránh chồng lấn với phần eBPF do
chính đồ án tự hiện thực ở tầng thu thập.

\subsection{containerd --- runtime chạy container}

\techlogo{containerd.png}{Logo containerd.}{fig:logo-containerd}
containerd là runtime container (\textit{container runtime}) cấp thấp, một dự án CNCF,
đảm nhiệm việc kéo ảnh, quản lý vòng đời và chạy container theo chuẩn OCI. Trong cụm, nó
nằm ngay dưới Kubernetes và là nơi thực sự khởi chạy các container của mọi pod, bao gồm
cả pod thu thập eBPF (vốn cần quyền đặc thù để truy cập tầng nhân). Đồ án dùng containerd
vì đây là runtime mặc định, gọn nhẹ và được Kubernetes hỗ trợ trực tiếp qua giao diện CRI,
không cần lớp trung gian như Docker Engine.

\subsection{NATS JetStream --- trục thông điệp xuất bản/đăng ký}

\techlogo{nats.png}{Logo NATS.}{fig:logo-nats}
NATS là một hệ thống truyền thông điệp (\textit{messaging}) hiệu năng cao theo mô hình
xuất bản--đăng ký (\textit{publish/subscribe}); JetStream là lớp bổ sung khả năng lưu trữ
bền vững và phát lại thông điệp. Trong kiến trúc của đồ án, NATS JetStream là \emph{trục
xương sống} nối tầng phân tích với các dịch vụ SOC: pipeline phát các cảnh báo lên các chủ
đề (\textit{subject}), còn aggregator, alert-bridge và incident-service đăng ký nhận. Cách
ghép nối lỏng này tách rời bên sinh sự kiện khỏi bên tiêu thụ, cho phép mở rộng hoặc thay
thế từng dịch vụ độc lập. NATS được ưu tiên hơn Kafka hay RabbitMQ vì nhẹ, độ trễ thấp và
chi phí vận hành nhỏ --- phù hợp với khối lượng cảnh báo của hệ thống mà không cần bộ máy
nặng của Kafka.

\subsection{PostgreSQL --- cơ sở dữ liệu quan hệ}

\techlogo{postgresql.png}{Logo PostgreSQL.}{fig:logo-postgres}
PostgreSQL là hệ quản trị cơ sở dữ liệu quan hệ mã nguồn mở, nổi tiếng về độ tin cậy và
tuân thủ chuẩn. Đồ án dùng PostgreSQL làm kho lưu trữ bền vững cho sự cố
(\textit{incident}), cảnh báo, hành động phản hồi và \emph{nhật ký kiểm toán}
(\textit{audit log}) --- những dữ liệu đòi hỏi tính toàn vẹn cao. Hai đặc tính quyết định
việc lựa chọn là: bảo đảm giao dịch ACID (cần thiết để nhật ký kiểm toán không bị mất hay
mâu thuẫn) và kiểu dữ liệu \texttt{JSONB} cho phép lưu linh hoạt các trường ngữ cảnh bán
cấu trúc (ví dụ thông tin pod đính kèm mỗi sự cố) mà vẫn truy vấn hiệu quả. So với các cơ
sở dữ liệu tài liệu như MongoDB, PostgreSQL giữ được ràng buộc quan hệ giữa sự cố -- cảnh
báo -- hành động; so với MySQL, nó hỗ trợ \texttt{JSONB} và chỉ mục biểu thức tốt hơn cho
nhu cầu này.

\subsection{Prometheus --- thu thập số liệu}

\techlogo{prometheus.png}{Logo Prometheus.}{fig:logo-prometheus}
Prometheus là hệ thu thập số liệu (\textit{metrics}) theo mô hình kéo (\textit{pull}) và
lưu trữ dưới dạng chuỗi thời gian, một dự án CNCF. Định kỳ nó truy vấn các điểm cuối
\texttt{/metrics} mà mỗi dịch vụ phơi ra và lưu lại theo thời gian. Trong đồ án, Prometheus
thu thập số liệu về tốc độ sự kiện, độ trễ xử lý và sức khỏe các dịch vụ SOC; chính những
số liệu này về sau được dùng làm dữ liệu \emph{đánh giá} định lượng ở Chương~\ref{ch:eval}.
Prometheus được chọn vì là chuẩn thực tế cho giám sát trong hệ sinh thái cloud-native, tích
hợp sẵn với Kubernetes và có mô hình truy vấn (PromQL) mạnh mà vẫn nhẹ.

\subsection{Grafana --- trực quan hóa và bảng điều khiển}

\techlogo{grafana.png}{Logo Grafana.}{fig:logo-grafana}
Grafana là công cụ trực quan hóa dữ liệu thành các bảng điều khiển (\textit{dashboard})
tương tác. Nó đọc số liệu từ Prometheus (và các nguồn khác) rồi hiển thị thành biểu đồ thời
gian thực cho người vận hành. Trong hệ thống, Grafana là mặt giám sát trực quan: chuyên
viên phân tích theo dõi tình hình cụm, tốc độ cảnh báo và sức khỏe dịch vụ qua các bảng do
Grafana dựng. Grafana được chọn vì ghép cặp tự nhiên với Prometheus, hỗ trợ nhiều loại biểu
đồ và phân quyền xem theo vai trò --- phù hợp với nhu cầu của một bảng điều khiển an ninh.

\subsection{Docker --- đóng gói ảnh container}

\techlogo{docker.png}{Logo Docker.}{fig:logo-docker}
Docker là công cụ đóng gói ứng dụng cùng toàn bộ phụ thuộc thành một \emph{ảnh}
(\textit{image}) bất biến theo chuẩn OCI, bảo đảm phần mềm chạy nhất quán giữa máy phát
triển và cụm thật. Trong đồ án, mỗi thành phần (collector, pipeline, sáu dịch vụ SOC) được
đóng gói thành một ảnh Docker riêng; tính bất biến của ảnh là nền tảng cho khả năng tái lập
của hệ thống. Docker được chọn vì là định dạng đóng gói phổ biến nhất, tương thích trực
tiếp với runtime \texttt{containerd} của cụm.

\subsection{Tekton --- tích hợp liên tục gốc Kubernetes}

\techlogo{tekton.png}{Logo Tekton.}{fig:logo-tekton}
Tekton là khung tích hợp liên tục (CI) \emph{gốc Kubernetes}: mỗi bước xây dựng, kiểm thử
và đóng gói ảnh được mô tả khai báo và chạy ngay trong cụm dưới dạng pod. Trong đồ án,
Tekton đảm nhiệm việc tự động xây và đóng gói ảnh OCI cho từng thành phần mỗi khi mã nguồn
thay đổi. Tekton được chọn thay cho các dịch vụ CI bên ngoài (như GitHub Actions) vì nó
chạy nội cụm, không phụ thuộc hạ tầng ngoài, giúp toàn bộ pipeline khép kín và kiểm toán
được --- phù hợp với một hệ thống an ninh đề cao tính tự chủ.

\subsection{Harbor --- kho ảnh container riêng}

\techlogo{harbor.png}{Logo Harbor.}{fig:logo-harbor}
Harbor là kho ảnh (\textit{registry}) container mã nguồn mở, một dự án CNCF, bổ sung các
tính năng an ninh chuỗi cung ứng so với một registry thông thường. Trong đồ án, Harbor là
nơi lưu trữ các ảnh do Tekton tạo ra; trước khi một ảnh được phép triển khai, Harbor quét
lỗ hổng và cho phép ký ảnh để bảo đảm tính toàn vẹn. Harbor được chọn thay cho một registry
công cộng (như Docker Hub) vì nó là kho \emph{riêng}, có kiểm soát truy cập và quét bảo mật
--- yêu cầu thiết yếu khi vận hành một hệ thống an ninh.

\subsection{ArgoCD --- triển khai theo mô hình GitOps}

\techlogo{argo.png}{Logo Argo (ArgoCD).}{fig:logo-argo}
ArgoCD là công cụ triển khai liên tục (CD) theo mô hình GitOps: trạng thái mong muốn của
cụm được khai báo trong một kho Git, và ArgoCD liên tục đồng bộ cụm thật về đúng khai báo
đó, tự dọn dẹp tài nguyên thừa và tự chữa lành khi có sai lệch. Trong đồ án, ArgoCD triển
khai toàn bộ hệ thống từ các manifest trong Git, nên mọi thay đổi đều có dấu vết và có thể
khôi phục. ArgoCD được chọn vì mô hình khai báo của nó làm cho hệ thống \emph{tái lập được}
và \emph{kiểm toán được} --- một yêu cầu phi chức năng quan trọng đã nêu ở
Chương~\ref{ch:overview}.

Ba công cụ Tekton, Harbor và ArgoCD tạo thành một chuỗi cung ứng phần mềm \emph{gốc
Kubernetes} khép kín trong cụm: Tekton xây ảnh, Harbor lưu và quét ảnh, ArgoCD triển khai
theo Git. Kết hợp với mô hình ảnh bất biến của Docker, chuỗi này là cơ sở kỹ thuật cho tính
tái lập và kiểm toán của toàn nền tảng.

\subsection{Kiến trúc nền tảng tổng thể và khác biệt với cách tiếp cận truyền thống}

Hình~\ref{fig:platform} tổng hợp toàn bộ các tầng nền tảng và đặt trên đó là ứng dụng
đặc thù của đồ án (bộ thu thập, pipeline, các dịch vụ SOC). Điểm cốt lõi mà sơ đồ muốn
làm nổi bật nằm ở mũi tên nhấn mạnh tầng nhân: ứng dụng quan sát lưu lượng \emph{ngay
tại tầng nhân} của mỗi nút, thay vì ở biên mạng hay bên trong mã ứng dụng.

\begin{figure}[H]
  \centering
  \resizebox{0.95\linewidth}{!}{% So do nen tang tong the (platform stack): tu nhan Linux (duoi) toi ung dung (tren).
% Nhan manh diem quan sat o tang nhan. Dung o Chuong 2.
\begin{tikzpicture}[font=\footnotesize, >={Stealth[length=2.2mm]},
    layer/.style={draw, rounded corners=2pt, minimum width=11.4cm,
                  minimum height=1.02cm, align=center},
  ]
  % Ung dung o TREN cung, nhan Linux o DUOI cung
  \node[layer, fill=red!12] (l6)
    {\textbf{Ứng dụng đặc thù} --- collector (DaemonSet) $\cdot$ pipeline đa-DSA $\cdot$ SOC (incident / response / web-console)};
  \node[layer, fill=orange!14, below=0.16cm of l6] (l5)
    {\textbf{Phân phối \& tự động hóa} --- GitOps ArgoCD $\cdot$ CI Tekton $\cdot$ Harbor registry};
  \node[layer, fill=green!12, below=0.16cm of l5] (l4)
    {\textbf{Dịch vụ nền} --- NATS JetStream $\cdot$ PostgreSQL $\cdot$ Prometheus / Grafana};
  \node[layer, fill=blue!12, below=0.16cm of l4] (l3)
    {\textbf{Điều phối \& mạng} --- Kubernetes $+$ CNI Calico (\texttt{NetworkPolicy})};
  \node[layer, fill=gray!10, below=0.16cm of l3] (l2)
    {\textbf{Runtime container} --- \texttt{containerd}};
  \node[layer, fill=gray!20, line width=1pt, below=0.16cm of l2] (l1)
    {\textbf{Phần cứng \& nhân Linux 6.8} --- eBPF (kprobe, CO-RE, BTF) $\cdot$ RingBuf $\cdot$ LPM-Trie};

  % Truc chieu cao tang
  \draw[->, gray!55, line width=0.8pt] ($(l1.south west)+(-0.35,0)$) -- ($(l6.north west)+(-0.35,0)$);
  \node[gray!70, rotate=90, font=\scriptsize, anchor=south] at ($(l1.south west)!0.5!(l6.north west)+(-0.55,0)$)
    {hạ tầng $\rightarrow$ ứng dụng};

  % Chu thich nhan manh tang nhan (node co text width -> khong tran le)
  \node[draw=red!70, line width=0.9pt, rounded corners, align=left, font=\scriptsize,
        fill=red!4, text width=3.05cm, right=0.55cm of l1] (note)
    {\textbf{Điểm quan sát:} ứng dụng đọc sự kiện kết nối \emph{ngay ở tầng nhân} --- không cần tác nhân trong ứng dụng, không phụ thuộc biên mạng.};
  \draw[->, red!70, line width=0.9pt] (note.west) -- (l1.east);
\end{tikzpicture}
}
  \caption{Sơ đồ nền tảng tổng thể: các tầng từ nhân Linux tới ứng dụng đặc thù.}
  \label{fig:platform}
\end{figure}

Sự khác biệt so với các kiến trúc truyền thống có thể nhìn trên bốn trục. Thứ nhất, về
\emph{vị trí quan sát}: an ninh chu vi truyền thống đặt thiết bị giám sát ở biên mạng
(hướng Bắc -- Nam) và gần như không thấy lưu lượng nội bộ giữa các dịch vụ, trong khi
nền tảng này quan sát toàn bộ lưu lượng Đông -- Tây ngay ở tầng nhân, đúng tinh thần
Zero-Trust \cite{nist800207}. Thứ hai, về \emph{cách thu thập}: thay vì cài tác nhân
vào từng ứng dụng (xâm lấn, phụ thuộc ngôn ngữ) hay phản chiếu gói toàn mạng (tốn kém),
eBPF cho phép thu thập không xâm lấn và độ phủ toàn nút chỉ với một \texttt{DaemonSet}.
Thứ ba, về \emph{mô hình vận hành}: hạ tầng truyền thống thường cấu hình thủ công và
trạng thái trôi dạt theo thời gian, còn nền tảng này khai báo bằng GitOps với ảnh bất
biến nên tái lập được và kiểm toán được. Thứ tư, về \emph{vòng phản hồi}: thay vì phân
tích theo lô rồi xử lý sự cố thủ công như các hệ SIEM cổ điển, hệ thống phát hiện trực
tuyến bằng cấu trúc dữ liệu luồng và khép kín vòng phản hồi bằng \texttt{NetworkPolicy}
có người duyệt. Chính bốn khác biệt này tạo nên giá trị của việc xây dựng ứng dụng an
ninh \emph{trên} một nền tảng cloud-native hiện đại, và là tiền đề cho phần thiết kế hệ
thống ở Chương~\ref{ch:design}.

\section{Kết luận chương}

Chương này đã thiết lập đầy đủ nền tảng lý thuyết và công nghệ cho phần còn lại của đồ
án. Bốn trụ lý thuyết được lựa chọn không phải ngẫu nhiên mà bám sát đúng đường đi của
một sự kiện trong hệ thống: sự kiện kết nối được \emph{thu thập} ở tầng nhân bằng eBPF
(Mục~\ref{sec:theory-ebpf}), được \emph{đếm tần suất} bằng Count-Min Sketch để phát hiện
quét cổng (Mục~\ref{sec:theory-cms}), rồi được \emph{phân tích quan hệ} bằng các giải
thuật đồ thị Tarjan và BFS, đồng thời \emph{phân loại} theo vùng địa chỉ bằng LPM-Trie
(Mục~\ref{sec:graph-ip}). Với mỗi giải thuật, chương đã trình bày ý tưởng, công thức,
mã giả và ví dụ chạy tay; riêng phần ước lượng tần suất chỉ nêu lại kết quả cô đọng, còn
toàn bộ chứng minh đầy đủ (giới hạn lỗi, cận dưới không gian, so sánh Misra--Gries và
Count Sketch) đã được phát triển chi tiết trong Báo cáo học phần tốt nghiệp (CNPM~1) và
được tham chiếu khi cần.

Bên cạnh các giải thuật, chương cũng đã giới thiệu từng công nghệ nền tảng --- từ
Kubernetes, Calico, containerd đến NATS, PostgreSQL, Prometheus/Grafana và chuỗi cung
ứng phần mềm --- cùng lý do lựa chọn, và kết lại bằng một bức tranh tổng thể về nền tảng
cùng bốn điểm khác biệt so với các kiến trúc truyền thống. Những cơ sở này là tiền đề
trực tiếp cho việc phân tích và thiết kế hệ thống ở Chương~\ref{ch:design}.
